% GENERATED FILE. Edit canonical lessons, then run npm run generate:books.
% canonical-source-hash: fnv1a64:4e7cbdd8c142c577
% canonical-lessons: ZH-C514-shandian, ZH-C514-linyu, ZH-C514-leiyu, ZH-C514-dongjiang, ZH-C514-ronghua

\chapter{When the Weather Turns}
\label{ch:zh-when-the-weather-turns}
\hlchaptermodality{514}

\begin{chapteropening}
\textbf{By the end of this chapter:} \emph{I can talk about wind and storms, getting caught in the rain, things that freeze and snow that melts.}

Complete the last lesson of chapter 514: I can talk about wind and storms, getting caught in the rain, things that freeze and snow that melts.
\end{chapteropening}

\section[shǎndiàn]{shǎndiàn — lightning}

\label{lesson:ZH-C514-shandian}



\begin{quote}
Before the new one: say the Mandarin for to struggle, then the Mandarin for with one's own hands.
\end{quote}

\subsection*{You'll want to know: shǎndiàn}
\emph{shǎndiàn} — \textquotedblleft{}lightning\textquotedblright{}.

\begin{quote}
\emph{wǒ zuì pà dǎléi hé shǎndiàn} — Thunder and lightning are what I'm most afraid of.
\end{quote}

Its characters come later; for now it is a word for the ear.

\begin{cousinweb}
\emph{shǎndiàn} is built from \emph{shǎn} and \emph{diàn} 'lightning', both inherited from Middle Chinese, \emph{syemX} and \emph{denH} (Guangyun, in Baxter's transcription). The dictionary glosses the first as the look of something coming out of a door, and explains the second as the flash when yin and yang clash; the sense 'to flash' is not given for the first.
\end{cousinweb}

\begin{grammarlens}[title={What you've built}]
One more word for saying how it was.
\end{grammarlens}

\subsection*{Your turn}
\begin{itemize}
\raggedright
  \item \emph{Say it:} \emph{shǎndiàn}
  \item \emph{Say it:} \emph{shǎndiàn}, once more
  \item \emph{Say it:} \emph{wǒ zuì pà dǎléi hé shǎndiàn} — Thunder and lightning are what I'm most afraid of.
  \item \emph{Recall:} say the Mandarin for to struggle, then the Mandarin for with one's own hands, then say \emph{shǎndiàn} again
\end{itemize}

\begin{tcolorbox}[breakable,colback=teal!4,colframe=teal!35!black,title={Before you move on}]
What does shǎndiàn mean? (\textquotedblleft{}Lightning\textquotedblright{}.) Say it once more.
\end{tcolorbox}
\section[línyǔ]{línyǔ — to get caught in the rain}

\label{lesson:ZH-C514-linyu}



\begin{quote}
Before the new one: say the Mandarin for with one's own hands, then the Mandarin for a flash of lightning.
\end{quote}

\subsection*{You'll want to know: línyǔ}
\emph{línyǔ} — \textquotedblleft{}to get caught in the rain\textquotedblright{}.

\begin{quote}
\emph{zuótiān línyǔ yǐhòu wǒ gǎnjué hěn lěng} — After I got caught in the rain yesterday, I felt cold.
\end{quote}

Its characters come later; for now it is a word for the ear.

\begin{cousinweb}
\emph{línyǔ} is built from \emph{lín} 'to pour on' and \zh{雨} 'rain', both inherited from Middle Chinese, \emph{lim} and \emph{hjuX} (Guangyun, in Baxter's transcription), which glosses the first as pouring water over something: to be poured on by the rain.
\end{cousinweb}

\begin{grammarlens}[title={What you've built}]
One more word for saying how it was.
\end{grammarlens}

\subsection*{Your turn}
\begin{itemize}
\raggedright
  \item \emph{Say it:} \emph{línyǔ}
  \item \emph{Say it:} \emph{línyǔ}, once more
  \item \emph{Say it:} \emph{zuótiān línyǔ yǐhòu wǒ gǎnjué hěn lěng} — After I got caught in the rain yesterday, I felt cold.
  \item \emph{Recall:} say the Mandarin for with one's own hands, then the Mandarin for a flash of lightning, then say \emph{línyǔ} again
\end{itemize}

\begin{tcolorbox}[breakable,colback=teal!4,colframe=teal!35!black,title={Before you move on}]
What does línyǔ mean? (\textquotedblleft{}To get caught in the rain\textquotedblright{}.) Say it once more.
\end{tcolorbox}
\section[léiyǔ]{léiyǔ — thunderstorm}

\label{lesson:ZH-C514-leiyu}



\begin{quote}
Before the new one: say the Mandarin for a flash of lightning, then the Mandarin for to get caught in the rain.
\end{quote}

\subsection*{You'll want to know: léiyǔ}
\emph{léiyǔ} — \textquotedblleft{}thunderstorm\textquotedblright{}.

\begin{quote}
\emph{wàipó shuō jīntiān wǎnshang yǒu léiyǔ} — Grandma says there'll be a thunderstorm tonight.
\end{quote}

Its characters come later; for now it is a word for the ear.

\begin{cousinweb}
\emph{léiyǔ} is built from \emph{léi} 'thunder' and \zh{雨} 'rain', both inherited from Middle Chinese, \emph{lwoj} and \emph{hjuX} (Guangyun, in Baxter's transcription), which explains the first, after an older book, as yin and yang striking each other.
\end{cousinweb}

\begin{grammarlens}[title={What you've built}]
One more word for saying how it was.
\end{grammarlens}

\subsection*{Your turn}
\begin{itemize}
\raggedright
  \item \emph{Say it:} \emph{léiyǔ}
  \item \emph{Say it:} \emph{léiyǔ}, once more
  \item \emph{Say it:} \emph{wàipó shuō jīntiān wǎnshang yǒu léiyǔ} — Grandma says there'll be a thunderstorm tonight.
  \item \emph{Recall:} say the Mandarin for a flash of lightning, then the Mandarin for to get caught in the rain, then say \emph{léiyǔ} again
\end{itemize}

\begin{tcolorbox}[breakable,colback=teal!4,colframe=teal!35!black,title={Before you move on}]
What does léiyǔ mean? (\textquotedblleft{}Thunderstorm\textquotedblright{}.) Say it once more.
\end{tcolorbox}
\section[dòngjiāng]{dòngjiāng — to be frozen stiff}

\label{lesson:ZH-C514-dongjiang}



\begin{quote}
Before the new one: say the Mandarin for to get caught in the rain, then the Mandarin for a thunderstorm.
\end{quote}

\subsection*{You'll want to know: dòngjiāng}
\emph{dòngjiāng} — \textquotedblleft{}to be frozen stiff\textquotedblright{}.

\begin{quote}
\emph{wǒ děng le yí ge xiǎoshí shǒu dōu dòngjiāng le} — I waited an hour, and my hands were frozen stiff.
\end{quote}

Its characters come later; for now it is a word for the ear.

\begin{cousinweb}
\emph{dòngjiāng} is built from \emph{dòng} 'to freeze' and \emph{jiāng} 'stiff', both inherited from Middle Chinese, \emph{tuwngH} and \emph{kjang} (Guangyun, in Baxter's transcription), which glosses the first 'ice, freezing' and the second as falling over flat.
\end{cousinweb}

\begin{grammarlens}[title={What you've built}]
One more word for saying how it was.
\end{grammarlens}

\subsection*{Your turn}
\begin{itemize}
\raggedright
  \item \emph{Say it:} \emph{dòngjiāng}
  \item \emph{Say it:} \emph{dòngjiāng}, once more
  \item \emph{Say it:} \emph{wǒ děng le yí ge xiǎoshí shǒu dōu dòngjiāng le} — I waited an hour, and my hands were frozen stiff.
  \item \emph{Recall:} say the Mandarin for to get caught in the rain, then the Mandarin for a thunderstorm, then say \emph{dòngjiāng} again
\end{itemize}

\begin{tcolorbox}[breakable,colback=teal!4,colframe=teal!35!black,title={Before you move on}]
What does dòngjiāng mean? (\textquotedblleft{}To be frozen stiff\textquotedblright{}.) Say it once more.
\end{tcolorbox}
\section[rónghuà]{rónghuà — to melt; to thaw}

\label{lesson:ZH-C514-ronghua}



\begin{quote}
Before the new one: say the Mandarin for a thunderstorm, then the Mandarin for to be frozen stiff.
\end{quote}

\subsection*{You'll want to know: rónghuà}
\emph{rónghuà} — \textquotedblleft{}to melt; to thaw\textquotedblright{}.

\begin{quote}
\emph{tiānqì nuǎnhuo le xuěrén dōu rónghuà le} — It turned warm, and the snowmen all melted.
\end{quote}

Its characters come later; for now it is a word for the ear.

\begin{cousinweb}
\emph{rónghuà} is built from \emph{róng} and \emph{huà} 'to change', both inherited from Middle Chinese, \emph{yuwng} and \emph{xwaeH} (Guangyun, in Baxter's transcription). The dictionary glosses the first 'harmonious' and 'bright' and, after an older book, as steam rising from a cooking pot; the second it glosses 'to change, to transform'. How the pair came to mean melting is usually explained as warmth changing a thing, but it is not settled.
\end{cousinweb}

\begin{grammarlens}[title={What you've built}]
One more word for saying how it was.
\end{grammarlens}

\subsection*{Your turn}
\begin{itemize}
\raggedright
  \item \emph{Say it:} \emph{rónghuà}
  \item \emph{Say it:} \emph{rónghuà}, once more
  \item \emph{Say it:} \emph{tiānqì nuǎnhuo le xuěrén dōu rónghuà le} — It turned warm, and the snowmen all melted.
  \item \emph{Recall:} say the Mandarin for a thunderstorm, then the Mandarin for to be frozen stiff, then say \emph{rónghuà} again
\end{itemize}

\begin{tcolorbox}[breakable,colback=teal!4,colframe=teal!35!black,title={Before you move on}]
What does rónghuà mean? (\textquotedblleft{}To melt; to thaw\textquotedblright{}.) Say it once more.
\end{tcolorbox}
